\subsection{Nagata rings}

DEFINITION 2. --- A ring A is said to be a Nagata ring if it is Noetherian and if, for every prime ideal p of A, the Noetherian integral domain A/p is Japanese (no. 1, def. 1).

Examples. --- 1) Every finitely generated algebra over a field is a Nagata ring (no. 1, example).

\begin{enumerate}
\def\labelenumi{\arabic{enumi})}
\setcounter{enumi}{1}
\item
  Every complete local Noetherian ring is a Nagata ring (no. 2, th. 2).
\item
  The ring Z is a Nagata ring (no. 1, example and corollary of prop. 1).
\item
  One can show (exerc. 30) that every finitely generated algebra over a Nagata ring is a Nagata ring.
\end{enumerate}

PROPOSITION 4. --- Let A be a Nagata ring.

\begin{enumerate}
\def\labelenumi{\alph{enumi})}
\item
  Every finite A-algebra is a Nagata ring.
\item
  For every multiplicative subset S of A, the ring \(\mathbf{S}^{-1}\mathbf{A}\) is a Nagata ring.
\item
  Let B be a finite A-algebra, \(\rho: A \to B\) the canonical homomorphism. For every prime ideal p of B, the ring B/p, which is a finite algebra over the Japanese ring \(A/\rho^{-1}(p)\) , is Japanese (no. 1, prop. 2).
\item
  Let q be a prime ideal of \(S^{-1}A\) ; then there exists a prime ideal p of A such that \(q = S^{-1}p\) . The ring ( \(S^{-1}A\) )/q is a ring of fractions of the Japanese ring A/p, hence is Japanese (no. 1, remark 2).
\end{enumerate}

THEOREM 3 (Zariski-Nagata). --- Let A be a semi-local Noetherian ring. The following conditions are equivalent:

\begin{enumerate}
\def\labelenumi{(\roman{enumi})}
\item
  A is a Nagata ring;
\item
  for every prime ideal \(\mathfrak{p}\) of A, the \(\kappa(\mathfrak{p})\) -algebra \(\kappa(\mathfrak{p}) \otimes_{\mathbf{A}} \hat{\mathbf{A}}\) is separable;
\item
  for every reduced A-algebra R, the ring \(R \otimes_{A} \hat{A}\) is reduced.
\end{enumerate}

Let us first prove the equivalence of conditions (ii) and (iii). The implication (iii) \(\Rightarrow\) (ii) is trivial; suppose conversely that A satisfies condition (ii). Then, for every A-algebra K which is a field, the ring K \(\otimes_{\mathbf{A}}\hat{\mathbf{A}}\) is reduced. Now let C be a reduced finitely generated A-algebra; the ring C, being Noetherian, is isomorphic to a subring of a finite product K \(_1\) \(\times\) \(\cdots\) \(\times\) K \(_n\) of fields (IV, § 2, no. 5, prop. 10); since \(\hat{\mathbf{A}}\) is flat over A, the ring C \(\otimes_{\mathbf{A}}\hat{\mathbf{A}}\) is isomorphic to a subring of the reduced ring \(\prod_{i}(\mathrm{K}_{i}\otimes_{\mathbf{A}}\hat{\mathbf{A}})\) , hence is reduced. Finally, let R be an arbitrary reduced A-algebra; then R is the union of the filtered family \((\mathrm{C}_{\alpha})\) of its finitely generated subalgebras, and \(R \otimes_{A} \hat{A}\) is the inductive limit of the filtered family \((\mathrm{C}_{\alpha} \otimes_{A} \hat{A})\) of reduced rings, hence is reduced.

Let us show that (ii) implies (i). Let p be a prime ideal of A; the field of fractions K of the ring A/p is identified with \(\kappa(p)\) , and the K-algebra \(K \otimes_{A/p} (\widehat{A/p})\) is identified with \(\kappa(p) \otimes_{A/p} \widehat{A}/p\widehat{A}\) , hence with \(\kappa(p) \otimes_{A} \widehat{A}\) . If \(\kappa(p) \otimes_{A} \widehat{A}\) is a \(\kappa(p)\) -separable algebra, the ring A/p is Japanese (no. 2, prop. 3).

Let us prove the implication (i) \(\Rightarrow\) (ii) by induction on dim(A). It is evident if dim(A) = 0, since then A is Artinian, hence complete. Let \(n\) be an integer \(>0\) ; consider the following hypothesis:

\((\mathbf{R}_n)\left\{
\begin{array}{ll}\text{for every local Noetherian Nagata ring C of dimension } < n\text{ and every prime}\\ \text{ideal } \mathfrak{r}\text{ of C, the ring }\kappa(\mathfrak{r})\otimes_{\mathrm{C}}\hat{\mathrm{C}}\text{ is reduced.} \end{array}
\right.\)

Let A be a semi-local Noetherian Nagata ring of dimension n, let p be a prime ideal of A, and let L be a finite extension of the field \(\kappa(\mathfrak{p})\); it suffices to prove, under hypothesis (R \(_{n}\) ), that the ring \(L \otimes_{A} \hat{A}\) is reduced. Let B denote the integral closure of A/p in L; since A/p is Japanese, B is a finite A-algebra, hence a semi-local Nagata ring (prop. 4). Let m \(_{1}\) , \ldots, m \(_{r}\) be the maximal ideals of B; the ring \(L \otimes_{A} \hat{A}\) is identified with a ring of fractions of \(B \otimes_{A} \hat{A}\), and the latter is identified with the product of the completions of the local rings B \(_{m_{i}}\) (no. 3, lemma 1). It therefore suffices to prove that, for every maximal ideal m of B, the ring B \(_{m}\) is reduced (II, § 2, no. 6, prop. 17). The ring B \(_{m}\) is local, integrally closed, and Nagata (prop. 4), and one has dim(B \(_{m}\) ) ≤ dim(B) ≤ dim(A) = n (VIII, § 1, no. 3, prop. 6 and § 2, no. 3, th. 1). Changing notations, one is reduced to proving, under hypothesis (R \(_{n}\) ), that for every local Noetherian integrally closed Nagata ring A of dimension ≤ n, the ring \(\hat{A}\) is reduced, that is to say (no. 3, lemma 2) that \(\hat{A}_{p'}\) is reduced for every prime ideal p′ ∈ Ass(\(\hat{A}\)). Since this is immediate if dim(A) = 0, we may assume dim(A) > 0. Let then x be a nonzero element of m \(_{A}\) and q' a prime ideal of \(\hat{A}\), minimal among those containing \(x\hat{A} + p'\); since \(\hat{A}_{p'}\) is identified with a ring of fractions of the ring \(\hat{A}_{q'}\), it suffices to prove that the latter is reduced (II, § 2, no. 6, prop. 17). By lemma 3, the ideal q' is associated to the \(\hat{A}\)-module \(\hat{A}/x\hat{A}\); since \(\hat{A}\) is flat over A, the inverse image q of q' in A is associated to the A-module A/xA (IV, § 2, no. 6, cor. 1 to th. 2). Since the ring A is assumed to be integrally closed, this implies that q has height 1 (VII, § 1, no. 6, prop. 10), hence that the ring \(A_{q}\) is a discrete valuation ring (loc. cit., no. 3, corollary to th. 2 and no. 6, th. 4). Since A/q is a Nagata ring of dimension < n, the ring \(\kappa(q) \otimes_{A/q} \hat{A}/q\hat{A}\) is reduced by hypothesis (R \(_{n}\) ). The ring \(\kappa(q) \otimes_{A} \hat{A}\), which is isomorphic to it, is reduced, and consequently so is the ring \(\kappa(q) \otimes_{A_{q}} \hat{A}_{q'}\), which is a ring of fractions of it. One may therefore apply Lemma 4 of Section 3 to the canonical homomorphism \(A_{q} \to \hat{A}_{q'}\), and we conclude that the ring \(\hat{A}_{q'}\) is reduced, which is what we wanted to prove. Theorem 3 is thus proved.

COROLLARY 1. --- The completion of a local and reduced Nagata ring is reduced.

It suffices indeed to set R = A in Thm. 3, (iii).

COROLLARY 2 (Chevalley). --- Let A be a reduced algebra of finite type over a field, and let p be a prime ideal of A. The completion of the local ring A \(_{p}\) is reduced.

Since A is reduced, the local ring \(A_{p}\) is reduced; moreover, A is a Nagata ring (Example 1), hence \(A_{p}\) is a Nagata ring (prop. 4), and cor. 2 follows from cor. 1, applied to the ring \(A_{p}\) .

COROLLARY 3. --- Let k be a field of characteristic 0, and let A be a local Noetherian k-algebra. For A to be a Nagata ring, it is necessary and sufficient that, for every prime ideal p of A, the ring (A/p) be reduced.

Indeed, since the fields \(\kappa(p)\) are of characteristic 0, it is equivalent to say that the algebras \(\kappa(p) \otimes_{A} \hat{A} = \kappa(p) \otimes_{A/p} (\widehat{A/p})\) are reduced or that they are separable (A, V, p.~117, th. 1), which shows that the stated condition is sufficient (th. 3, (ii) \(\Rightarrow\) (i)); it is moreover necessary (th. 3, (i) \(\Rightarrow\) (iii) with \(R = A/p\) ).

